% ============================================================================
% Chapitre 8 — Déploiement et mise en production
% ============================================================================
\chapter{Déploiement et mise en production}

\section*{Introduction}

Ce chapitre décrit les procédures d'installation et de configuration de
\caimp{} dans les différents contextes d'utilisation : développement local,
serveur \textsc{vps} de petite taille, et déploiement de production avec
haute disponibilité.

\section{Modes de déploiement}

\caimp{} supporte trois modes de déploiement, adaptés à des niveaux de ressources
et d'exigences différents :

\begin{table}[H]
\centering
\caption{Modes de déploiement et ressources requises}
\label{tab:deploy-modes}
\begin{tabularx}{\textwidth}{lXXX}
\toprule
\textbf{Mode} & \textbf{Commande} & \textbf{RAM min.} & \textbf{Usage} \\
\midrule
Complet    & \texttt{make up}                & 8~Go   & Production, développement \\
Low-RAM    & \texttt{bash scripts/start-test.sh} & 4~Go & VPS économique, CI \\
Production & Docker Swarm / Kubernetes       & 16~Go+ & Haute disponibilité \\
\bottomrule
\end{tabularx}
\end{table}

\section{Procédure d'installation complète}

\subsection{Prérequis}

\begin{itemize}
  \item Docker Engine 24+ et Docker Compose v2.20+ ;
  \item 20~Go d'espace disque libre ;
  \item Accès Internet pour le premier téléchargement des images et des modèles Ollama.
\end{itemize}

\subsection{Étapes d'installation}

\begin{enumerate}[leftmargin=2.2cm]
  \item \textbf{Cloner le dépôt} :
\begin{lstlisting}[language=YAML]
git clone https://github.com/your-org/caimp.git && cd caimp
\end{lstlisting}

  \item \textbf{Configurer les variables d'environnement} :
\begin{lstlisting}[language=YAML]
cp .env.example .env
# Éditer .env : renseigner les mots de passe et JWT_SECRET
# JWT_SECRET : openssl rand -hex 64
\end{lstlisting}

  \item \textbf{Démarrer la plateforme} :
\begin{lstlisting}[language=YAML]
make up               # Build + démarrage de tous les services
make ps               # Vérifier que tout est "healthy"
\end{lstlisting}

  \item \textbf{Télécharger les modèles LLM} (première exécution uniquement) :
\begin{lstlisting}[language=YAML]
make pull-models      # llama3.1:8b-instruct-q4_K_M + nomic-embed-text
\end{lstlisting}

  \item \textbf{Accéder au tableau de bord} :
    Ouvrir \texttt{http://localhost} dans un navigateur.\\
    Identifiants par défaut : \texttt{admin@caimp.local} / \texttt{Admin1234!}

  \item \textbf{Enrôler un serveur cible} : générer un jeton d'enrôlement via
    l'\api{} Admin, puis l'exécuter sur le serveur à surveiller.
\end{enumerate}

\subsection{Configuration des variables d'environnement}

\begin{table}[H]
\centering
\caption{Variables d'environnement obligatoires}
\label{tab:env-vars}
\small
\begin{tabularx}{\textwidth}{lXl}
\toprule
\textbf{Variable} & \textbf{Description} & \textbf{Obligatoire} \\
\midrule
\texttt{POSTGRES\_PASSWORD}    & Mot de passe PostgreSQL        & Oui \\
\texttt{REDIS\_PASSWORD}       & Mot de passe Redis             & Oui \\
\texttt{NATS\_PASSWORD}        & Mot de passe NATS              & Oui \\
\texttt{JWT\_SECRET}           & Clé secrète JWT (64 hex chars) & Oui \\
\texttt{CA\_KEY\_PASSPHRASE}   & Passphrase CA interne          & Oui \\
\texttt{GRAFANA\_ADMIN\_PASSWORD} & Mot de passe Grafana       & Oui \\
\texttt{OLLAMA\_LLM\_MODEL}    & Modèle LLM (défaut: llama3.1:8b) & Non \\
\texttt{SPLUNK\_HOST}          & Hôte Splunk (activer l'intégration) & Non \\
\texttt{GITHUB\_WEBHOOK\_SECRET} & Secret webhook GitHub       & Non \\
\texttt{LOG\_LEVEL}            & DEBUG/INFO/WARNING (défaut: INFO) & Non \\
\bottomrule
\end{tabularx}
\end{table}

\section{Mode Low-RAM (4 Go)}

Pour les machines contraintes en mémoire (4~Go de RAM), le script
\texttt{start-test.sh} applique le fichier \texttt{docker-compose.test.yml} :

\begin{itemize}
  \item Réduction des limites mémoire de tous les services ;
  \item Exclusion d'Ollama du démarrage par défaut
    (\texttt{--profile full} requis) ;
  \item Remplacement de Llama 3.1 8B par Llama 3.2 1B (1.3~Go) ;
  \item Concurrence AI Worker réduite à 1 ;
  \item Création automatique d'un fichier swap de 4~Go.
\end{itemize}

Consommation mémoire mesurée en mode test : \textbf{2.3~Go} (pile de base sans Ollama).

\section{Sauvegarde et restauration}

\begin{lstlisting}[language=YAML, caption={Procédure de sauvegarde PostgreSQL}]
# Sauvegarde complète
docker compose exec postgres \
  pg_dump -U caimp -d caimp -F c -f /tmp/caimp_backup.dump

docker cp caimp-postgres:/tmp/caimp_backup.dump ./backups/

# Restauration
docker cp ./backups/caimp_backup.dump caimp-postgres:/tmp/
docker compose exec postgres \
  pg_restore -U caimp -d caimp -c /tmp/caimp_backup.dump
\end{lstlisting}

\section{Monitoring de la plateforme elle-même}

Prometheus scrape les endpoints \texttt{/metrics} de chaque service toutes les
15 secondes. Les métriques clés surveillées sont :

\begin{itemize}
  \item \textbf{Telemetry Writer} : \texttt{tw\_metrics\_ingested\_total},
    \texttt{tw\_anomalies\_detected\_total}, latence p99 ;
  \item \textbf{AI Worker} : \texttt{ai\_explanations\_generated\_total},
    \texttt{ai\_ollama\_latency\_seconds} ;
  \item \textbf{Alert Engine} : \texttt{alert\_engine\_anomalies\_suppressed\_total}
    (déduplication) ;
  \item \textbf{Query API} : latence HTTP p99, taux d'erreur 5xx.
\end{itemize}

Ces métriques sont visualisées dans Grafana (port 3001), accessible avec les
identifiants configurés dans \texttt{GRAFANA\_ADMIN\_PASSWORD}.

\section{Génération de métriques de démonstration}

Pour alimenter le tableau de bord avec des données réalistes sans disposer d'un
vrai serveur, \caimp{} fournit un agent de démonstration léger
(\texttt{services/agent/demo\_agent.py}) qui lit les métriques réelles
(\textsc{cpu}, \textsc{ram}, disque) de la machine hôte et les transmet à la
plateforme toutes les 20 secondes.

\paragraph{Prérequis}
\begin{lstlisting}[language=YAML]
pip install psutil requests
\end{lstlisting}

\paragraph{Utilisation de base — un serveur virtuel avec métriques réelles}
\begin{lstlisting}[language=YAML]
python services/agent/demo_agent.py
\end{lstlisting}

\noindent Cela enregistre le serveur virtuel \texttt{demo-server-01} dans
l'organisation par défaut et commence immédiatement à envoyer des métriques.

\paragraph{Simulation d'alertes critiques — trois terminaux simultanés}

\begin{lstlisting}[language=YAML, caption={Terminal 1 — métriques réelles de la machine}]
python services/agent/demo_agent.py
\end{lstlisting}

\begin{lstlisting}[language=YAML, caption={Terminal 2 — db-primary avec RAM critique}]
CAIMP_SERVER_ID=00000000-0000-0000-0000-000000000013 \
CAIMP_SERVER_NAME=db-primary \
CAIMP_RAM=97 \
python services/agent/demo_agent.py
\end{lstlisting}

\begin{lstlisting}[language=YAML, caption={Terminal 3 — worker-01 avec CPU critique}]
CAIMP_SERVER_ID=00000000-0000-0000-0000-000000000014 \
CAIMP_SERVER_NAME=worker-01 \
CAIMP_CPU=99 \
python services/agent/demo_agent.py
\end{lstlisting}

\noindent En moins de deux minutes, l'AI Worker génère des explications pour
chaque dépassement de seuil critique, visibles dans le tableau de bord et
interrogeables via l'assistant \textsc{ia} (chatbot).

\begin{table}[H]
\centering
\caption{Variables d'environnement de l'agent de démonstration}
\label{tab:demo-agent-vars}
\small
\begin{tabularx}{\textwidth}{lXl}
\toprule
\textbf{Variable} & \textbf{Description} & \textbf{Valeur par défaut} \\
\midrule
\texttt{CAIMP\_SERVER\_ID}   & UUID du serveur virtuel            & \texttt{...0011} \\
\texttt{CAIMP\_SERVER\_NAME} & Nom affiché dans le tableau de bord & \texttt{demo-server-01} \\
\texttt{CAIMP\_CPU}          & Forcer le CPU\,\% (0--100)         & valeur réelle \\
\texttt{CAIMP\_RAM}          & Forcer la RAM\,\% (0--100)         & valeur réelle \\
\texttt{CAIMP\_DISK}         & Forcer le Disque\,\% (0--100)      & valeur réelle \\
\texttt{CAIMP\_INTERVAL}     & Intervalle d'envoi (secondes)      & \texttt{20} \\
\texttt{CAIMP\_TW\_URL}      & URL du telemetry-writer            & valeur interne \\
\bottomrule
\end{tabularx}
\end{table}

\section{Surveillance d'une machine virtuelle réelle}

Pour surveiller un vrai serveur Linux (avec systemd), il faut générer un jeton
d'enrôlement via l'\textsc{api} Admin, copier l'agent sur la machine cible,
puis exécuter le programme d'installation fourni.

\paragraph{Étape 1 — Générer un jeton d'enrôlement}
Depuis la machine faisant tourner \caimp{} :
\begin{lstlisting}[language=YAML]
curl -s -X POST http://localhost/api/v1/admin/servers/enroll \
  -H "Authorization: Bearer <admin_jwt>" \
  -H "Content-Type: application/json" \
  -d '{"name":"ma-vm","org_id":"<org_uuid>"}' | python3 -m json.tool
\end{lstlisting}
La réponse contient un \texttt{enrollment\_token} à usage unique.

\paragraph{Étape 2 — Copier l'agent sur la VM cible}
\begin{lstlisting}[language=YAML]
scp -r services/agent/ utilisateur@<ip-vm>:~/caimp-agent/
\end{lstlisting}

\paragraph{Étape 3 — Exécuter le programme d'installation sur la VM cible}
\begin{lstlisting}[language=YAML]
ssh utilisateur@<ip-vm>
cd ~/caimp-agent
sudo bash install.sh \
  --token <enrollment_token> \
  --url http://<ip-hote-caimp>
\end{lstlisting}

\noindent Le programme d'installation crée un utilisateur système
\texttt{caimp-agent}, installe un environnement virtuel Python, enrôle l'agent
(génération d'une paire de clés cryptographiques et soumission d'un \textsc{csr}),
et démarre le service systemd \texttt{caimp-agent} automatiquement au démarrage.

\paragraph{Étape 4 — Vérifier que l'agent fonctionne}
\begin{lstlisting}[language=YAML]
systemctl status caimp-agent
journalctl -u caimp-agent -f
\end{lstlisting}

\noindent Le serveur apparaît dans le tableau de bord \caimp{} dans les
30~secondes suivant la réception du premier lot de métriques.

\section{Accès aux tableaux de bord}

\caimp{} s'exécute entièrement sur la machine où \texttt{docker compose up} a
été lancé — le poste de l'administrateur en développement, ou un serveur dédié
en production. Toutes les interfaces web sont servies par le proxy inverse nginx
sur le port~80.

\begin{table}[H]
\centering
\caption{URLs des interfaces web}
\label{tab:web-urls}
\begin{tabularx}{\textwidth}{lXl}
\toprule
\textbf{Interface} & \textbf{Description} & \textbf{URL} \\
\midrule
Tableau de bord \caimp{} & Portail admin — connexion, serveurs, chatbot \textsc{ia}
  & \texttt{http://localhost} \\
Grafana & Graphiques, séries temporelles et historique des alertes
  & \texttt{http://localhost:3001} \\
\bottomrule
\end{tabularx}
\end{table}

\paragraph{Accès depuis une autre machine}
Si \caimp{} tourne sur un serveur d'adresse IP \texttt{192.168.1.10}, remplacer
\texttt{localhost} par cette adresse dans chaque URL :
\begin{lstlisting}[language=YAML]
http://192.168.1.10        # Tableau de bord CAIMP
http://192.168.1.10:3001   # Grafana
\end{lstlisting}

\paragraph{Identifiants par défaut}
\begin{itemize}
  \item \textbf{Tableau de bord \caimp{}} : \texttt{admin@caimp.local} /
    \texttt{Admin1234!}
  \item \textbf{Grafana} : \texttt{admin} / valeur de
    \texttt{GRAFANA\_ADMIN\_PASSWORD} dans \texttt{.env}
\end{itemize}

\paragraph{Développement vs.\ production}
En développement, la plateforme tourne sur la même machine que le navigateur —
aucune configuration réseau supplémentaire n'est nécessaire. En production, la
plateforme doit tourner sur un serveur accessible par toutes les machines
surveillées ; l'installateur de l'agent doit recevoir l'adresse IP publique ou
privée du serveur via l'option \texttt{--url}.

\section*{Conclusion}

La procédure d'installation en six étapes rend \caimp{} accessible en moins
de dix minutes, même sur une machine à ressources limitées. La présence d'un
mode low-RAM et de scripts de sauvegarde automatisés garantit l'opérabilité
dans des environnements contraints. Le chapitre suivant analyse les résultats
obtenus et discute les perspectives d'évolution.
